\lysh{34325}{4}

\begin{alist}
\item Το τριώνυμο $x^2 - x + (\lambda - \lambda^2)$ έχει $\alpha = 1, \beta = -1, \gamma = \lambda - \lambda^2$ και διακρίνουσα
\[\Delta = \beta^2 - 4\alpha\gamma = (-1)^2 - 4 \cdot 1 \cdot (\lambda - \lambda^2) =\]
$= 1 - 4\lambda + 4\lambda^2 = (1 - 2\lambda)^2 \ge 0$, για κάθε $\lambda \in \mathbb{R}$.
Επειδή $\Delta \ge 0$, για κάθε $\lambda \in \mathbb{R}$, η εξίσωση $x^2 - x + \lambda - \lambda^2 = 0$ έχει πραγματικές ρίζες.

\item Η εξίσωση έχει δύο άνισες πραγματικές ρίζες αν και μόνο αν:
\begin{align*}
\Delta > 0 \Leftrightarrow (1 - 2\lambda)^2 > 0 \ &\Leftrightarrow \\
\Leftrightarrow 1 - 2\lambda &\ne 0 \Leftrightarrow \lambda \ne \dfrac{1}{2}.
\end{align*}

\item Οι ρίζες της εξίσωσης (1) είναι οι:
\[x_{1,2} = \dfrac{-(-1) \pm \sqrt{(1 - 2\lambda)^2}}{2} = \dfrac{1 \pm (1 - 2\lambda)}{2} = \begin{cases} \dfrac{1 + 1 - 2\lambda}{2} = 1 - \lambda \\ \dfrac{1 - 1 + 2\lambda}{2} = \lambda \end{cases}\]
Έχουμε ότι:
\begin{align*}
0 < d(x_1, x_2) < 2 \ &\Leftrightarrow \\
0 < |x_1 - x_2| < 2 \ &\Leftrightarrow \\
0 < |1 - \lambda - \lambda| < 2 \ &\Leftrightarrow \\
0 < |1 - 2\lambda| < 2 \ &\Leftrightarrow
\end{align*}
Δηλαδή,
\begin{align*}
0 < |1 - 2\lambda| \ &\Leftrightarrow \\
1 - 2\lambda \ne 0 \ &\Leftrightarrow \\
2\lambda \ne 1 \ &\Leftrightarrow \\
\lambda &\ne \dfrac{1}{2}
\end{align*}
και
\begin{align*}
|1 - 2\lambda| < 2 \ &\Leftrightarrow \\
-2 < 1 - 2\lambda < 2 \ &\Leftrightarrow \\
-3 < -2\lambda < 1 \ &\Leftrightarrow \\
-\dfrac{1}{2} &< \lambda < \dfrac{3}{2}.
\end{align*}
Οπότε, τελικά
\[\lambda \in \left(-\dfrac{1}{2}, \dfrac{1}{2}\right) \cup \left(\dfrac{1}{2}, \dfrac{3}{2}\right).\]
\end{alist}
